% !TEX root = ../supporting_information.tex

\SISection{Uncertainty quantification of
\texorpdfstring{\coTwo{}}{CO2} migration and containment}
\label{sec:si-migration-containment-uq}

\SISubsection{Quantities of interest and outcome classes}

Primary containment responses quantify whether \coTwo{} enters the top seal
and the corresponding mass fraction. The primary migration response is the
maximum upward distance reached by \coTwo{} relative to the injection
interval. Secondary responses include the fractions of injected \coTwo{}
retained in the storage reservoir, entering the modeled fault and overburden,
and partitioned between dissolved and free phases. All mass fractions use a
common injected-mass denominator and are evaluated at prespecified output
times.

Outcome classes distinguish consistent containment, localized fault entry,
substantial upward migration, and top-seal or overburden interaction. Their
spatial masks and numerical thresholds are defined before inspecting the
complete ensemble so that classification does not adapt to the observed
results. \todo{Provide the final storage-reservoir, fault, top-seal, and
overburden masks; injected-mass denominator; free and dissolved component
accounting; maximum-migration calculation; numerical entry tolerances;
evaluation times; and class thresholds.}

\SISubsection{Design-based geologic effects}

The common 12 independent realizations per geology form a balanced nested
design. For continuous response $Q_{g,r}$, geologic effects are estimated from
the four prescribed factors and selected interactions with realizations nested
within geology. Effects are interpreted over the uniform scenario design, not
as sensitivities weighted by field occurrence probabilities. Entry quantities
with a point mass at zero are separated into an occurrence model and a
conditional-magnitude model. \todo{State the final model form, contrasts,
standardization, uncertainty intervals, and interaction-selection rule.}

For each geology, the 12 stochastic realizations also estimate the conditional
median, robust spread, empirical outcome frequencies, and empirical CDFs. The
medoid and low/high benchmark cases are excluded from these estimators. This
separation prevents deterministic diagnostic fields from being counted as
random draws from the conditional fault-property distribution.

\SISubsection{Between- and within-geology contributions}

Let $\bar Q_g$ and $s_g^2$ denote the sample mean and variance among the 12
independent realizations for geology $g$. With $G=162$, the average
within-geology contribution is estimated as
\begin{equation}
  \widehat V_{\mathrm{within}}
  =\frac{1}{G}\sum_{g=1}^{G}s_g^2.
\end{equation}
The variance among $\bar Q_g$ includes finite-replicate noise, so the
between-geology design contribution is estimated as
\begin{equation}
  \widehat V_{\mathrm{between}}
  =\max\left[0,
  \operatorname{Var}_g(\bar Q_g)
  -\frac{1}{G}\sum_{g=1}^{G}\frac{s_g^2}{12}\right].
\end{equation}
Normalized contributions divide each quantity by their sum. Robust spreads
and outcome-class disagreement accompany this variance summary when responses
are non-Gaussian, zero-inflated, or multimodal.

\SISubsection{Targeted enrichment and confirmatory analysis}

Phase-1 results nominate geologies using prespecified criteria: more than one
outcome class, large robust conditional spread, medoid--ensemble disagreement,
stress benchmarks on opposite sides of a physical threshold, and coverage of
the geologic design space. Twelve geologies receive 20 new independent
realizations. The six that remain most realization sensitive receive 40
further realizations. Selection statistics use only earlier simulations;
newly added simulations evaluate whether the apparent sensitivity persists.
\todo{Report the exact thresholds, selected geology identifiers, normalized
geologic-space coverage, and all selection decisions.}

\SISubsection{Representative and stress benchmark analyses}

For each geology, representative-fault performance is quantified by the
difference between the medoid result and the stochastic median, the medoid's
position in the empirical stochastic CDF, and agreement with the majority
outcome class. Low/high benchmarks are compared as coherent state stresses but
are excluded from empirical CDFs, variance contributions, and outcome
probabilities. \todo{Add threshold-sensitivity analyses and complete benchmark
results for all quantities of interest.}
